\documentclass[11pt]{article}
% ===========================================================================
%  Shared preamble for the AP Calculus AB exam-preparation worksheets.
%  Input right after \documentclass. Files beginning with "_" are not compiled
%  on their own. Navy and gold, to match the course's lesson pages.
% ===========================================================================
\usepackage[letterpaper,top=2.3cm,bottom=2.2cm,left=1.7cm,right=1.7cm,headheight=28pt,footskip=30pt]{geometry}
\usepackage{amsmath,amssymb}
\usepackage[T1]{fontenc}
\usepackage[utf8]{inputenc}
\usepackage{lmodern}
\usepackage{xcolor}
\usepackage[most]{tcolorbox}
\usepackage{enumitem}
\usepackage{fancyhdr}
\usepackage{array}
\usepackage{tikz}
\usepackage{pgfplots}
\pgfplotsset{compat=1.18}

\definecolor{navy}{HTML}{161E54}
\definecolor{ink}{HTML}{27357E}
\definecolor{gold}{HTML}{B98A00}
\definecolor{goldbg}{HTML}{FFFDF3}
\definecolor{inkbg}{HTML}{F6F7FD}
\definecolor{rule}{HTML}{C9CFF0}
\definecolor{okgreen}{HTML}{1B7A44}
\definecolor{okbg}{HTML}{F1F8F3}

\pagestyle{fancy}
\fancyhf{}
\fancyhead[L]{\footnotesize Name: \underline{\hspace{4cm}}}
\fancyhead[C]{\textbf{\color{navy}Integration Academy}\\[-2pt]{\scriptsize Calculus AB \textperiodcentered{} AP\textsuperscript{\textregistered} Exam Preparation}}
\fancyfoot[L]{\tiny\color{black!55}\parbox[t]{0.78\textwidth}{AP\textsuperscript{\textregistered} is a registered trademark of the College Board, which was not involved in the production of, and does not endorse, this material. \textcopyright{} Integration Academy.}}
\fancyfoot[R]{\scriptsize\thepage}
\renewcommand{\headrulewidth}{1.2pt}
\renewcommand{\headrule}{\hbox to\headwidth{\color{gold}\leaders\hrule height \headrulewidth\hfill}}
\renewcommand{\footrulewidth}{0pt}
\setlength{\parindent}{0pt}
\setlength{\parskip}{4pt}

% Title block: \wstitle{code}{title}{unit line}
\newcommand{\wstitle}[3]{\begin{center}{\small\bfseries\color{gold}WORKSHEET #1}\\[3pt]{\LARGE\bfseries\color{navy}#2}\\[3pt]{\small\color{black!60}#3}\end{center}\vspace{2pt}}
% Section heading
\newcommand{\sech}[1]{\par\vspace{10pt}{\large\bfseries\color{navy}#1}\par\nobreak\vspace{2pt}{\color{rule}\hrule height0.8pt}\vspace{7pt}}
% A level tag printed beside a group of questions
\newcommand{\level}[1]{\par\vspace{6pt}{\footnotesize\bfseries\color{gold}\MakeUppercase{#1}}\par\nobreak\vspace{2pt}}

\newtcolorbox{ideabox}{colframe=ink,colback=inkbg,coltitle=white,colbacktitle=ink,fonttitle=\bfseries\small,title={Key ideas},arc=2pt,boxrule=0.8pt,left=8pt,right=8pt}
\newtcolorbox{example}[1]{colframe=rule,colback=white,coltitle=navy,colbacktitle=inkbg,fonttitle=\bfseries\small,title={#1},arc=2pt,boxrule=0.8pt,borderline west={3pt}{0pt}{ink},left=10pt,right=8pt}
\newtcolorbox{frbox}[1]{colframe=gold,colback=goldbg,coltitle=white,colbacktitle=gold,fonttitle=\bfseries\small,title={#1},arc=2pt,boxrule=0.8pt,breakable,left=8pt,right=8pt}
\newtcolorbox{solbox}{colframe=okgreen,colback=okbg,coltitle=white,colbacktitle=okgreen,fonttitle=\bfseries\small,title={Worked solutions},arc=2pt,boxrule=0.8pt,breakable,left=8pt,right=8pt}
\newtcolorbox{scorebox}[1]{colframe=navy,colback=white,coltitle=white,colbacktitle=navy,fonttitle=\bfseries\small,title={#1},arc=2pt,boxrule=0.8pt,breakable,left=8pt,right=8pt}

\newcommand{\soln}{\par\smallskip\textbf{\color{ink}Solution.}\ }
\newcommand{\result}[1]{\par\smallskip\textbf{#1}}
% blank working space of a given height
\newcommand{\work}[1]{\par\vspace{#1}}
% a point value in the margin of a scoring guide
\newcommand{\pts}[1]{\hfill{\footnotesize\bfseries\color{gold}[#1]}}

% Coordinate grid for graph questions: \begin{apgraph}{xmin}{xmax}{ymin}{ymax} ... \end{apgraph}
\newenvironment{apgraph}[4]{\begin{center}\begin{tikzpicture}\begin{axis}[
  axis lines=middle, xmin=#1, xmax=#2, ymin=#3, ymax=#4,
  grid=both, grid style={black!12}, xtick distance=1, ytick distance=1,
  ticklabel style={font=\scriptsize}, xlabel={$x$}, ylabel={$y$},
  width=10cm, height=7.2cm, samples=80, clip=false,
  every axis plot/.append style={line width=1.1pt, color=ink}]}{\end{axis}\end{tikzpicture}\end{center}}
\newcommand{\closeddot}[2]{\addplot[only marks, mark=*, mark size=2.2pt, color=ink] coordinates {(#1,#2)};}
\newcommand{\opendot}[2]{\addplot[only marks, mark=*, mark size=2.2pt, color=ink, mark options={fill=white}] coordinates {(#1,#2)};}

% Multiple-choice options: four across, or stacked when the options are long.
\newcommand{\opts}[4]{\par\vspace{3pt}\begin{tabular}{@{}p{0.235\linewidth}p{0.235\linewidth}p{0.235\linewidth}p{0.235\linewidth}@{}}(A)\ #1 & (B)\ #2 & (C)\ #3 & (D)\ #4\end{tabular}\par\vspace{9pt}}
\newcommand{\optsv}[4]{\par\vspace{3pt}\begin{tabular}{@{}l@{\ \ }p{0.9\linewidth}@{}}(A) & #1\\(B) & #2\\(C) & #3\\(D) & #4\end{tabular}\par\vspace{9pt}}

\fancyhead[R]{\footnotesize\color{black!60}Worksheet 6.12}
\begin{document}

\wstitle{6.12}{Integrals from Graphs and Tables}{Unit 6: Integration and Accumulation of Change}

\begin{ideabox}
\begin{itemize}[leftmargin=*,itemsep=2pt,topsep=2pt]
  \item From a graph, a definite integral is a sum of \textbf{signed areas} of triangles, rectangles and parts of circles.
  \item From a table, an integral can only be \textbf{approximated}, with a Riemann or trapezoidal sum.
  \item But $\displaystyle\int_{a}^{b}f'(x)\,dx=f(b)-f(a)$ is \textbf{exact}, from two values of $f$.
  \item The properties of integrals apply as usual: constants, sums and reversed limits.
\end{itemize}
\end{ideabox}

\sech{Worked examples}
\begin{example}{Example 1 \textperiodcentered{} Areas}
The graph of $g$ is the line segment from $(0,3)$ to $(6,-3)$. Evaluate $\displaystyle\int_{0}^{6}g(x)\,dx$ and $\displaystyle\int_{0}^{6}|g(x)|\,dx$.\soln Triangles of area $\dfrac92$ above the axis and $\dfrac92$ below.\result{$0$ and $9$.}
\end{example}

\begin{example}{Example 2 \textperiodcentered{} A table}
Selected values of $h$ are shown.\begin{center}\renewcommand{\arraystretch}{1.25}\begin{tabular}{|c|c|c|c|c|}\hline $x$ & $0$ & $2$ & $6$ & $7$ \\ \hline $h(x)$ & $5$ & $9$ & $3$ & $1$ \\ \hline\end{tabular}\end{center}Approximate $\displaystyle\int_{0}^{7}h(x)\,dx$ with a trapezoidal sum.\soln $\dfrac12(5+9)(2)+\dfrac12(9+3)(4)+\dfrac12(3+1)(1)=14+24+2$.\result{$40$.}
\end{example}

\begin{example}{Example 3 \textperiodcentered{} The integral of a derivative}
For the function $h$ above, assumed differentiable, find $\displaystyle\int_{0}^{7}h'(x)\,dx$.\soln $h(7)-h(0)$.\result{$-4$, exactly.}
\end{example}

\newpage
\sech{Practice}
Questions 1 to 7 use this graph of $f$, which consists of four line segments.
\begin{apgraph}{-4}{6}{-3}{5}
\addplot[color=ink, mark=none, line width=1.1pt] coordinates {(-3,0) (-1,4) (1,0) (3,-2) (5,0)};
\closeddot{-3}{0}\closeddot{5}{0}
\node[font=\small, color=ink] at (axis cs:0.6,3.2) {$y=f(x)$};
\end{apgraph}
\level{Core}
\begin{enumerate}[leftmargin=*,itemsep=2pt]
  \item Find $\displaystyle\int_{-3}^{1}f(x)\,dx$.\work{1.1cm}
  \item Find $\displaystyle\int_{1}^{5}f(x)\,dx$.\work{1.1cm}
  \item Find $\displaystyle\int_{-3}^{5}f(x)\,dx$.\work{1.1cm}
  \item Find $\displaystyle\int_{-3}^{5}|f(x)|\,dx$.\work{1.1cm}
\end{enumerate}

\level{Exam level}
\begin{enumerate}[leftmargin=*,itemsep=2pt,start=5]
  \item Find $\displaystyle\int_{-1}^{3}f(x)\,dx$.\work{2.1cm}
  \item Find $\displaystyle\int_{-3}^{1}\big[2f(x)+1\big]\,dx$.\work{2.1cm}
  \item Find $\displaystyle\int_{-1}^{3}f'(x)\,dx$.\work{2.1cm}
  \item A differentiable function $k$ has $k(2)=11$ and $\displaystyle\int_{2}^{9}k'(x)\,dx=-6$. Find $k(9)$.\work{2.1cm}
\end{enumerate}

\begin{frbox}{Question 9 \textperiodcentered{} Free-response practice}
The function $h$ is differentiable on $[0,10]$. Selected values are shown.\begin{center}\renewcommand{\arraystretch}{1.25}\begin{tabular}{|c|c|c|c|c|}\hline $t$ & $0$ & $4$ & $6$ & $10$ \\ \hline $h(t)$ & $8$ & $14$ & $11$ & $3$ \\ \hline\end{tabular}\end{center}
\begin{enumerate}[label=(\alph*),leftmargin=*,itemsep=2pt]
  \item Approximate $\displaystyle\int_{0}^{10}h(t)\,dt$ with a trapezoidal sum using the three subintervals of the table.\work{1.9cm}
  \item Find the exact value of $\displaystyle\int_{0}^{10}h'(t)\,dt$.\work{1.9cm}
  \item Find the exact value of $\displaystyle\int_{4}^{10}\big[h'(t)+2\big]\,dt$.\work{1.9cm}
\end{enumerate}
\end{frbox}

\newpage
\begin{solbox}
\begin{enumerate}[leftmargin=*,itemsep=4pt]
  \item A triangle of base 4 and height 4: $8$.
  \item A triangle of base 4 and depth 2: $-4$.
  \item $8-4=4$.
  \item $8+4=12$.
  \item On $[-1,1]$: $\dfrac12(2)(4)=4$. On $[1,3]$: $-\dfrac12(2)(2)=-2$. Value $2$.
  \item $2(8)+1(4)=20$.
  \item $f(3)-f(-1)=-2-4=-6$.
  \item $11+(-6)=5$.
  \item \begin{enumerate}[label=(\alph*),leftmargin=*,itemsep=3pt]
    \item $\dfrac12(8+14)(4)+\dfrac12(14+11)(2)+\dfrac12(11+3)(4)=44+25+28=97$.
    \item $h(10)-h(0)=-5$.
    \item $\big(h(10)-h(4)\big)+2(6)=-11+12=1$.
  \end{enumerate}
\end{enumerate}
\end{solbox}

\end{document}
